% Separate development, 9 September 2026.
% Extends the product reader of centro_electronico_registros.tex.
% Does not identify the spatial half-turn with a physical conjugation.
\subsubsection{Bidirectionality of the product reader and quotient memory}
\label{bim09:seccion}

The product reading preserves four arithmetic coordinates before
determining a signature. For $x=(i,j)\in D_9^2$, $D_9=\{1,\ldots,9\}$,
define
\[
 N(x)=ij,\qquad r(x)=1+((ij-1)\bmod9),\qquad
 q(x)=\frac{N(x)-r(x)}9,
\]
and the discrete reader
\[
 \ell(1,2,4,8,7,5)=(0,1,2,3,4,5),\qquad
 \ell(3)=\ell(6)=\ell(9)=0.
\]
Thus $N=r+9q$ simultaneously preserves the integer product, positive residue,
and arithmetic quotient. The function $f_\ell(x)=\ell(r(x))$ is the
coordinate used by the regional signature. The functions $N,r,q$, and
$f_\ell$ have a positional domain; the transport direction remains
in the oriented seed.

\paragraph{Advance, half-turn, and reading before transport.}

Let $\mathscr X_\times=D_9^2\times\{N,E,S,O\}$. The cardinal vectors
are $\nu_N=(-1,0)$, $\nu_E=(0,1)$, $\nu_S=(1,0)$,
$\nu_O=(0,-1)$. Positional addition is reduced to positive representatives
modulo nine. Define
\[
 P(x,d)=(x+\nu_d,d),\qquad H(x,d)=(x,\bar d),
 \qquad \nu_{\bar d}=-\nu_d.
\]
These operators satisfy $P^9=I$, $H^2=I$, and $HPH=P^{-1}$.

In the first fifty-four TPK transitions, the product channel
is activated in phases $2,5,8$ of each nonadic window. Each activation
reads the position before advance. After transitions $27$ and $54$,
$H$ is applied. Thus the eighteen positions read are, for a
seed $s$,
\[
 P^0s,P^1s,\ldots,P^8s,\quad
 HP^9s,\,P H P^9s,\ldots,P^8H P^9s,
\]
with the understanding that the positional reading omits only direction.
For any positional function $f$, put $f_n=f(P^ns)$, with
indices modulo nine. The word of raw readings is
\begin{equation}
 \mathcal R_f(s)=
 (f_0,f_1,\ldots,f_8,\ f_0,f_8,f_7,\ldots,f_1).
 \label{bim09:palabra}
\end{equation}
The six window sums, each comprising three activations, are
\begin{align}
 A_j^f(s)&=f_{3j}+f_{3j+1}+f_{3j+2},&&0\leq j<3,\label{bim09:ida}\\
 A_{3+j}^f(s)&=f_{-3j}+f_{-3j-1}+f_{-3j-2},&&0\leq j<3.
 \label{bim09:vuelta}
\end{align}
For $f=f_\ell$, the signature is
$E_{\rm sig}(s)=([A_0^f]_{10},\ldots,[A_5^f]_{10})$.
The decimal quotient of each accumulator is
$c_j^{10}=\lfloor A_j^f/10\rfloor$, and it preserves
$A_j^f=[A_j^f]_{10}+10c_j^{10}$.

\begin{proposition}[Half-turn and permutation of half-cycles]
\label{bim09:media-vuelta}
For every oriented seed and every positional function $f$,
\[
 \mathcal R_f(Hs)=\operatorname{rot}_9\mathcal R_f(s),\qquad
 A_j^f(Hs)=A_{j+3\bmod6}^f(s).
\]
The same permutation transports the decimal residues and quotients
of the accumulators. It applies, in particular, to $N,r,q$, and $f_\ell$.
\end{proposition}
\begin{proof}
The relation $P^nH=HP^{-n}$ and equality $f(Hs)=f(s)$ give
$f(P^nHs)=f(P^{-n}s)=f_{-n}$. Substitution into
\eqref{bim09:palabra} exchanges its two groups of nine entries.
Summation in groups of three produces the window identity. Euclidean
division by ten is performed componentwise and respects the
permutation.
\end{proof}

Advance also has an exact law on the accumulators:
\begin{align}
 A_j^f(Ps)-A_j^f(s)&=f_{3j+3}-f_{3j},&&0\leq j<3,
 \label{bim09:avance-ida}\\
 A_{3+j}^f(Ps)-A_{3+j}^f(s)&=f_{1-3j}-f_{-2-3j},&&0\leq j<3.
 \label{bim09:avance-vuelta}
\end{align}
The proof consists in canceling the two common summands of each window.
Transport of a signature therefore requires the boundary evaluations
appearing in these differences, in addition to the sums already
expressed.

\paragraph{Chronological inversion and boundary terms.}

Let $\gamma=(x_0,\ldots,x_n)$ be a recorded trajectory, and let
$\gamma^\dagger=(x_n,\ldots,x_0)$ be its chronological inversion, with each
edge oriented in the opposite direction. For a positional function $f$,
the reading before each advance is
\[
 C_f(\gamma)=\sum_{t=0}^{n-1}f(x_t).
\]
Then
\begin{equation}
 C_f(\gamma^\dagger)=C_f(\gamma)+f(x_n)-f(x_0).
 \label{bim09:frontera-inversion}
\end{equation}
If $n=3m$ and $A_j^f$ groups three consecutive edges,
\begin{equation}
 A_j^f(\gamma^\dagger)=A_{m-1-j}^f(\gamma)
 +f(x_{3(m-j)})-f(x_{3(m-j-1)}).
 \label{bim09:frontera-ventana}
\end{equation}
Both formulas follow because the reverse trajectory reads the arrival
endpoints of the original edges. The interior summands coincide;
the difference preserves the endpoints.

For an antisymmetric edge observable $a(x,y)=-a(y,x)$, its word
does transform according to
\[
 \mathcal R_a(\gamma^\dagger)
 =-\operatorname{rev}\mathcal R_a(\gamma).
\]
This identity acts on oriented edges. Identities
\eqref{bim09:frontera-inversion}--\eqref{bim09:frontera-ventana}
act on positional readings. The half-turn $H$ changes the
initial direction and preserves calendar order; chronological
inversion reverses the complete record. Their domains and reading rules
are specified separately.

\paragraph{Occupation cocycle and memory preserved on outward and return paths.}

For a cardinal edge $x\to y$, with positive representatives
$\widetilde x,\widetilde y$, the covering crossing is
\[
 b(x,y)=\frac{\widetilde x+\nu_d-\widetilde y}{9}\in\mathbb Z^2.
\]
The identity $b(y,x)=-b(x,y)$ holds. In the free category of recorded
paths, define the additive reader
\begin{equation}
 \mathfrak m(\gamma)=
 \left(\sum_{t=0}^{n-1}b(x_t,x_{t+1}),\ C_q(\gamma),\ n\right).
 \label{bim09:memoria-aditiva}
\end{equation}
Concatenation satisfies
$\mathfrak m(\gamma_2\circ\gamma_1)=
\mathfrak m(\gamma_1)+\mathfrak m(\gamma_2)$.
It is an additive cocycle on recorded trajectories: the record preserves
outward and return edges as different events.

By \eqref{bim09:frontera-inversion},
\begin{equation}
 \mathfrak m(\gamma^\dagger\circ\gamma)=
 \bigl(0,\ 2C_q(\gamma)+q(x_n)-q(x_0),\ 2n\bigr),
 \label{bim09:memoria-retorno}
\end{equation}
where $\gamma^\dagger\circ\gamma$ denotes the temporal concatenation
of the outward path followed by the return. The oriented sum of crossings cancels;
arithmetic occupation and length remain. This reader does not
descend to the quotient that identifies a path followed by its inverse
with the empty trajectory, since that quotient eliminates the events
preserved by the second and third coordinates.

There is also a purely boundary-based even reading:
\[
 V_b(\gamma)=\sum_{t=0}^{n-1}
 b(x_t,x_{t+1})b(x_t,x_{t+1})^{\mathsf T},\qquad
 V_b(\gamma^\dagger\circ\gamma)=2V_b(\gamma).
\]
The odd and even readings record distinct properties of
the same edges; both are computed before the path is forgotten.

\paragraph{Exact application to the signature fiber \texorpdfstring{$555555$}{555555}.}

The finite enumeration of $\mathscr X_\times$ contains $324$ seeds,
$26$ signatures, and $24$ seeds with signature $555555$. The stable
partition under $P,H$ refines this fiber into three classes of eight elements.
They can be labeled by their signature after one advance:
\[
 \mathscr C_{177},\quad\mathscr C_{717},\quad\mathscr C_{771}.
\]
The half-turn action is
\[
 H\mathscr C_{177}=\mathscr C_{177},\qquad
 H\mathscr C_{717}=\mathscr C_{771},\qquad
 H\mathscr C_{771}=\mathscr C_{717}.
\]
The attached finite certificate verifies these identities on all
representatives. In the twenty-four seeds, the six raw accumulators
of $f_\ell$ are exactly five. Hence the decimal quotients
$c_j^{10}$ vanish on this fiber. The additional information
is observed in the positions, integer evaluations, and arithmetic
quotients, which have readings different from that signature.

The fixed transverse coordinate $a$ of the cardinal path is $4$ or
$5$ in these seeds. During one turn of nine advances, each
value of the moving coordinate $u\in D_9$ is traversed once. Since
$\gcd(a,9)=1$, the positive residues $r(au)$ permute $D_9$.
This gives, without using a physical constant,
\[
 \sum_{u=1}^9au=45a,\qquad
 \sum_{u=1}^9r(au)=45,\qquad
 \sum_{u=1}^9q(au)=5(a-1).
\]
The subsequent return traverses the same positions in the opposite direction.
Over the eighteen activations,
\begin{equation}
 \sum N=90a,\qquad\sum r=90,\qquad
 \sum q=10(a-1)=
 \begin{cases}30,&a=4,\\40,&a=5.\end{cases}
 \label{bim09:ocupacion-30-40}
\end{equation}
The accumulated integer product is, respectively, $360$ or $450$.
The net winding is zero; the sum of quotients and the eighteen
incidences preserve the history of the reading. Each of the three
signature classes contains representatives of both transverse values.
The same signature and the same transport class of that finite quotient
can therefore coexist with distinct arithmetic occupations.

These results connect preserved memory with concrete
operations: permutation of half-cycles, boundary differences,
arithmetic quotients, and incidence records. The physical-mathematical
identification of any subsequent reader preserves its own
maps; the identities proved here belong to the internal dynamics
of the product channel.
